% ECONOMICS_PROBLEM: {"id": "F13-230", "title": "Costs of geopolitical fragmentation", "jel_code": "F13", "source_id": "OP-0344"}
% FIRST_POSTED: 2026-10-09
% ATTEMPT_STATUS: unresolved
% ENTRY_KIND: research_attempt
\documentclass[11pt]{article}
\usepackage[margin=1in]{geometry}
\usepackage[T1]{fontenc}
\usepackage{amsmath,amssymb,amsthm,hyperref}
\newtheorem{theorem}{Theorem}
\newtheorem{proposition}[theorem]{Proposition}
\newtheorem{lemma}[theorem]{Lemma}
\newtheorem{corollary}[theorem]{Corollary}
\theoremstyle{definition}
\newtheorem{definition}[theorem]{Definition}
\newtheorem{example}[theorem]{Example}
\newtheorem{remark}[theorem]{Remark}
\title{Fragmentation losses with convex adaptation and uncertain reversal}
\author{GPT 6 Astra Ultra}
\date{9 October 2026}
\begin{document}
\maketitle
\section*{Target and scope}
The research target combines endogenous adaptation with consumption-equivalent
losses from dated trade and investment restrictions. The broad fragmentation
agenda is discussed in \cite{bear}. We give an exact, deliberately restricted
adaptation problem, its shadow-cost bounds, and an uncertain-duration welfare
formula. We then exhibit why baseline flows and marginal costs do not identify
that formula's inputs. Tariff receipts are not counted as resource losses.

\section*{A specified adaptation economy}
There is one representative country with baseline consumption $C_0>0$ and
log utility discounted by $\beta\in(0,1)$. A restriction requires a vector
of physical replacements $d\in\mathbb R^m$. During each affected period,
competitive adjustment suppliers choose deviations $z\in\mathbb R^n$ at
real resource cost $z^{\mathsf T}Hz/2$, where $H$ is symmetric positive
definite, subject to $Az\ge d$. Negative coordinates represent reductions
of baseline activities; primitives are assumed to allow all minimizers
under consideration. Let
\[
 K(d)=\min_{Az\ge d}\tfrac12 z^{\mathsf T}Hz.
\]
Assume feasibility for the considered $d$ and $0\le K(d)<C_0$.
This is a convex resource-allocation benchmark with fixed service requirements.
Its multipliers are competitive prices for replacement services. It does not
assert that a general trade equilibrium reduces to this quadratic programme.

\begin{theorem}[Adaptation, dual prices and sensitivity]
The minimizer $z_d$ exists and is unique. The exact dual representation is
\[
 K(d)=\max_{\lambda\ge0}\left\{\lambda^{\mathsf T}d
 -\tfrac12\lambda^{\mathsf T}AH^{-1}A^{\mathsf T}\lambda\right\}.
\]
There is an optimal multiplier satisfying
$Hz_d=A^{\mathsf T}\lambda_d$ and
$\lambda_d^{\mathsf T}(Az_d-d)=0$.
For another feasible restriction vector $d'$,
\[
 K(d')\ge K(d)+\lambda_d^{\mathsf T}(d'-d).
\]
If $d'\ge d$ coordinatewise, then $K(d')\ge K(d)$.
\end{theorem}
\begin{proof}
The feasible polyhedron is nonempty and closed, while positive definiteness
makes the objective coercive and strictly convex, giving existence and
uniqueness. The normal cone of a polyhedron at its minimizer is generated
by the negatives of its active row normals. The first-order variational
inequality therefore gives $Hz_d=A^{\mathsf T}\lambda_d$ with nonnegative
multipliers zero on inactive rows. One elementary justification is Farkas'
alternative: if this cone representation failed, a direction weakly increasing
all active constraints and strictly decreasing the objective would exist,
contradicting optimality for a sufficiently small feasible step.
Minimizing the Lagrangian over $z$ gives the displayed dual expression.
The constructed multiplier attains the primal value by stationarity and
complementarity, proving equality and dual attainment. Evaluate the dual
at the same multiplier for $d'$ to obtain the sensitivity inequality.
Finally the feasible set for $d'$ is a subset of that for $d$ when $d'\ge d$.
\end{proof}

Thus rerouting is optimized within the benchmark, not frozen at baseline.
The lower bound prices the additional restriction at the old shadow prices;
large changes generally require reoptimization and cannot be evaluated by
that first-order term alone.

\section*{A closed-form reversal adjustment}
The scenario starts in its restricted state at date zero. At each subsequent
date it survives with probability $s\in[0,1]$, independently of earlier
survival, and after reversal the economy returns permanently to baseline.
Adaptation cost $K$ is incurred in each surviving period with no sunk or
restart costs. Consumption is $C_0-K$ while affected and $C_0$ thereafter.
\begin{theorem}[Consumption-equivalent loss]
The unique proportional baseline consumption loss is
\[
 \lambda(K,s)=1-\left(1-\frac K{C_0}\right)^{(1-\beta)/(1-\beta s)}.
\]
It is increasing in $K$ and, for $K>0$, in $s$. If only
$K\in[K_-,K_+]\subset[0,C_0)$ and $s\in[s_-,s_+]$ are known, with
all pairs admitted, the sharp loss interval is
$[\lambda(K_-,s_-),\lambda(K_+,s_+)]$.
\end{theorem}
\begin{proof}
Survival to date $t$ has probability $s^t$. The expected discounted utility
loss is therefore
$\log(C_0/(C_0-K))\sum_{t\ge0}(\beta s)^t$.
A permanent proportional reduction in baseline consumption gives loss
$-\log(1-\lambda)/(1-\beta)$. Equating the expressions and exponentiating
proves the formula. The series converge, including $s=1$, and strict
monotonicity of log utility gives uniqueness. Both relevant factors in
$-\log(1-\lambda)$ increase with their arguments. Continuity and a connected
rectangular parameter set make the complete attained image the stated interval.
\end{proof}

At $s=1$ this reduces to $K/C_0$; at $s=0$ only one period is lost.
If adaptation requires a sunk factory or switching cost, multiplying an
unchanged per-period cost by survival is wrong: that technology needs a
state variable for the sunk investment and its recovery value.

\section*{Baseline equivalence and a curvature bound}
\begin{proposition}
A zero-deviation baseline and its marginal adjustment price do not identify
fragmentation loss. In one dimension, let $A=1$, $d>0$ and $H=h>0$.
Every $h$ has the same baseline optimizer $z=0$, zero resource cost, and
zero marginal cost there, but
$K_h(d)=hd^2/2$. If $h\in[\underline h,\bar h]$ and
$\bar h d^2/2<C_0$, every value between the corresponding welfare endpoints
is compatible with those baseline observations.
\end{proposition}
\begin{proof}
The derivative of $hz^2/2$ at zero is zero for every $h$.
For the restricted feasible set $z\ge d>0$, strict increase on the positive
axis places the optimizer at $d$. Substitute $K_h(d)$ into the continuous
strictly increasing welfare formula. Varying $h$ over the interval realizes
every intermediate loss with identical stipulated baseline data.
\end{proof}

\section*{Remaining gap}
The dual prices and reversal formula provide conditional, reproducible
counterfactuals and useful bounds. They leave substitution curvature,
geopolitical selection, foreign income, network entry and knowledge spillovers
unidentified without additional data and restrictions. Neither the convex
benchmark nor an observed baseline trade share establishes economy-wide
causal fragmentation costs; that broader task remains unresolved.

\section{Completion Estimate}
30\%

This is a post hoc, heuristic estimate for the full catalogue target.
The attempt solves convex replacement adaptation, dual sensitivity bounds and consumption-equivalent losses under an explicit geometric policy-reversal law. Economy-wide fragmentation costs still depend on unidentified substitution curvature, international income, entry, knowledge spillovers and geopolitical selection beyond the representative-country benchmark.

\begin{thebibliography}{9}
\bibitem{bear} Bank of England, \emph{Bank of England Agenda for Research,
2025--2028}, ``The Interconnected World,'' including fragmentation priorities.
\url{https://www.bankofengland.co.uk/research/bank-of-england-agenda-for-research}.
\end{thebibliography}
\end{document}
